\chapter{Time, Crypto, OS, and Everything Else}\label{ch:time-crypto-os}

\index{time module}
\index{crypto module}
\index{os module}
\index{math module}
\index{regex module}

This chapter is a grand tour of the remaining standard library modules---the
functions that do not fit neatly into files, formats, or networking, but that
you will reach for constantly in real programs. We will cover timestamps and
date arithmetic, cryptographic hashing and encoding, UUID generation,
operating-system interaction, mathematics, and regular expressions. Each
section is self-contained, so feel free to jump to whichever module you need
right now and come back to the others later.

% ---------------------------------------------------------------------------
\section{The time Module}\label{sec:time-module}
% ---------------------------------------------------------------------------

\index{time}
\index{sleep}
\index{time\_format}
\index{time\_parse}
\index{datetime}

Time in Lattice is represented as a single integer: the number of
milliseconds since the Unix epoch (January~1, 1970, 00:00:00 UTC). This is
the same convention used by JavaScript's \texttt{Date.now()}, making it
straightforward to exchange timestamps with web APIs.

\subsection{Getting the Current Time}\label{sec:time-now}

\begin{lstlisting}[language=Lattice, caption={Getting the current timestamp}]
let now = time()
print(now)  // e.g., 1703456789012

// Measure execution time
let start = time()
// ... some computation ...
let elapsed = time() - start
print("Took " + to_string(elapsed) + " ms")
\end{lstlisting}

\lstinline|time()| returns the current wall-clock time as an \lstinline|Int|
in milliseconds. Under the hood, it calls \texttt{clock\_gettime} with
\texttt{CLOCK\_REALTIME} (see \filename{src/time\_ops.c}).

\subsection{Sleeping}\label{sec:sleep}

\begin{lstlisting}[language=Lattice, caption={Pausing execution}]
print("Starting...")
sleep(1000)  // Sleep for 1 second
print("Done!")

// Short delay for polling
sleep(100)  // 100 milliseconds
\end{lstlisting}

\lstinline|sleep| takes a duration in milliseconds. It uses
\texttt{nanosleep} for precise timing and handles \texttt{EINTR} gracefully.

\begin{warnbox}{sleep Is Not Available in WASM}
In the browser (Emscripten) build, \lstinline|sleep| raises a runtime error
because it would block the single-threaded event loop. If you need
time-delayed behavior in WASM, use callbacks or an async pattern instead.
\end{warnbox}

\subsection{Formatting Timestamps}\label{sec:time-format}

\begin{lstlisting}[language=Lattice, caption={Formatting a timestamp as a human-readable string}]
let now = time()

let readable = time_format(now, "%Y-%m-%d %H:%M:%S")
print(readable)  // "2024-12-25 14:30:45"

let date_only = time_format(now, "%B %d, %Y")
print(date_only)  // "December 25, 2024"

let time_only = time_format(now, "%I:%M %p")
print(time_only)  // "02:30 PM"
\end{lstlisting}

\lstinline|time_format| takes a timestamp (in epoch milliseconds) and a
format string using the standard \texttt{strftime} specifiers:

\begin{center}
\begin{tabular}{llll}
\toprule
\textbf{Specifier} & \textbf{Meaning} & \textbf{Specifier} & \textbf{Meaning} \\
\midrule
\texttt{\%Y} & 4-digit year & \texttt{\%H} & Hour (00--23) \\
\texttt{\%m} & Month (01--12) & \texttt{\%I} & Hour (01--12) \\
\texttt{\%d} & Day (01--31) & \texttt{\%M} & Minute (00--59) \\
\texttt{\%B} & Full month name & \texttt{\%S} & Second (00--59) \\
\texttt{\%b} & Abbreviated month & \texttt{\%p} & AM/PM \\
\texttt{\%A} & Full weekday name & \texttt{\%Z} & Timezone abbreviation \\
\texttt{\%a} & Abbreviated weekday & \texttt{\%\%} & Literal \% \\
\bottomrule
\end{tabular}
\end{center}

\subsection{Parsing Timestamps}\label{sec:time-parse}

\begin{lstlisting}[language=Lattice, caption={Parsing a date string into a timestamp}]
let timestamp = time_parse("2024-12-25 14:30:00", "%Y-%m-%d %H:%M:%S")
print(timestamp)  // epoch milliseconds

let formatted = time_format(timestamp, "%B %d, %Y")
print(formatted)  // "December 25, 2024"
\end{lstlisting}

\lstinline|time_parse| is the inverse of \lstinline|time_format|. It takes a
date string and a format pattern, and returns the corresponding epoch
timestamp in milliseconds. The format specifiers are the same as for
\lstinline|time_format|.

\begin{notebox}{Local Time}
Both \lstinline|time_format| and \lstinline|time_parse| work in the
system's local timezone by default. The underlying C functions
(\texttt{localtime\_r} and \texttt{mktime}) handle timezone conversions
and daylight saving time adjustments automatically.
\end{notebox}

\subsection{Extracting Date Components}\label{sec:time-components}

\index{time\_year}
\index{time\_month}
\index{time\_day}
\index{time\_hour}
\index{time\_minute}
\index{time\_second}
\index{time\_weekday}

For quick access to individual components without constructing a format
string, Lattice provides dedicated extraction functions:

\begin{lstlisting}[language=Lattice, caption={Extracting date and time components}]
let now = time()

print(time_year(now))     // 2024
print(time_month(now))    // 12 (1-based)
print(time_day(now))      // 25
print(time_hour(now))     // 14
print(time_minute(now))   // 30
print(time_second(now))   // 45
print(time_weekday(now))  // 3 (0=Sunday, 1=Monday, ..., 6=Saturday)
\end{lstlisting}

\subsection{Date Arithmetic}\label{sec:date-arithmetic}

\index{time\_add}

Since timestamps are just integers in milliseconds, date arithmetic is
straightforward:

\begin{lstlisting}[language=Lattice, caption={Date arithmetic with timestamps}]
let now = time()

// Add 7 days
let next_week = time_add(now, 7 * 24 * 60 * 60 * 1000)
print("Next week: " + time_format(next_week, "%Y-%m-%d"))

// Subtract 30 days
let last_month = time_add(now, -30 * 24 * 60 * 60 * 1000)
print("30 days ago: " + time_format(last_month, "%Y-%m-%d"))

// Useful constants (in milliseconds)
let SECOND = 1000
let MINUTE = 60 * SECOND
let HOUR = 60 * MINUTE
let DAY = 24 * HOUR

let deadline = time_add(now, 3 * HOUR + 30 * MINUTE)
print("Deadline: " + time_format(deadline, "%H:%M:%S"))
\end{lstlisting}

\lstinline|time_add| simply adds a delta in milliseconds to a timestamp.
You can also use plain addition (\lstinline|now + delta|), but
\lstinline|time_add| makes the intent clearer.

\subsection{The datetime Functions}\label{sec:datetime}

\index{datetime\_now}
\index{datetime\_from\_iso}
\index{datetime\_to\_iso}
\index{datetime\_from\_epoch}
\index{datetime\_to\_epoch}
\index{datetime\_format}
\index{datetime\_add\_duration}
\index{datetime\_sub}
\index{datetime\_to\_utc}
\index{datetime\_to\_local}
\index{timezone\_offset}

For more structured date/time work---especially when dealing with ISO\,8601
strings and UTC conversions---Lattice provides a family of
\lstinline|datetime_*| functions. These work with epoch \emph{seconds} (not
milliseconds) and return Maps with structured fields.

\begin{lstlisting}[language=Lattice, caption={Working with ISO 8601 dates}]
// Get current datetime as a Map
let now = datetime_now()
print(now)  // Map with year, month, day, hour, minute, second, etc.

// Parse an ISO 8601 string
let dt = datetime_from_iso("2024-12-25T14:30:00Z")
print(dt)

// Convert back to ISO 8601
let iso_str = datetime_to_iso(dt)
print(iso_str)  // "2024-12-25T14:30:00Z"

// Format with a custom pattern
let formatted = datetime_format(dt, "%B %d, %Y at %H:%M")
print(formatted)  // "December 25, 2024 at 14:30"
\end{lstlisting}

\begin{lstlisting}[language=Lattice, caption={Duration arithmetic with datetime}]
let start = datetime_from_iso("2024-01-15T09:00:00Z")

// Add a duration (in seconds)
let meeting_end = datetime_add_duration(start, 90 * 60)  // 90 minutes

// Compute the difference between two datetimes (in seconds)
let end = datetime_from_iso("2024-01-15T10:30:00Z")
let diff = datetime_sub(end, start)
print("Duration: " + to_string(diff / 60) + " minutes")  // 90 minutes
\end{lstlisting}

The \lstinline|datetime_to_utc| and \lstinline|datetime_to_local| functions
convert between UTC and local time. \lstinline|timezone_offset()| returns the
local timezone offset from UTC in seconds.

\begin{tipbox}{Which Time API to Use?}
For simple timing and logging, the \lstinline|time()| / \lstinline|time_format()|
functions are sufficient. For applications that exchange dates with APIs,
store them in databases, or perform calendar arithmetic across timezones,
use the \lstinline|datetime_*| family with ISO\,8601 strings.
\end{tipbox}

% ---------------------------------------------------------------------------
\section{The crypto Module}\label{sec:crypto-module}
% ---------------------------------------------------------------------------

\index{sha256}
\index{sha512}
\index{md5}
\index{hmac\_sha256}
\index{base64\_encode}
\index{base64\_decode}
\index{random\_bytes}

Lattice provides cryptographic primitives for hashing, message
authentication, encoding, and random number generation. When built with
OpenSSL support (\texttt{LATTICE\_HAS\_TLS}), these functions use OpenSSL's
optimized implementations. Without OpenSSL, Lattice falls back to pure-C
implementations that produce identical results---just more slowly.

\subsection{Hashing: sha256, sha512, md5}\label{sec:hashing}

\begin{lstlisting}[language=Lattice, caption={Computing cryptographic hashes}]
let message = "Hello, Lattice!"

let hash = sha256(message)
print(hash)  // 64-character hex string
// e.g., "a1b2c3d4e5f6..."

let long_hash = sha512(message)
print(long_hash)  // 128-character hex string

let legacy_hash = md5(message)
print(legacy_hash)  // 32-character hex string
\end{lstlisting}

All hash functions take a \lstinline|String| and return a lowercase
hexadecimal string. \lstinline|sha256| produces a 256-bit (64 hex character)
digest. \lstinline|sha512| produces a 512-bit (128 hex character) digest.
\lstinline|md5| produces a 128-bit (32 hex character) digest.

\begin{warnbox}{MD5 Is Not Secure}
MD5 is cryptographically broken and should never be used for security
purposes (password hashing, digital signatures, integrity verification of
untrusted data). It is included for backward compatibility with legacy
systems. Use \lstinline|sha256| or \lstinline|sha512| for new applications.
\end{warnbox}

\begin{lstlisting}[language=Lattice, caption={File integrity checking with sha256}]
fn file_checksum(path: String) -> String {
    let contents = read_file(path)
    return sha256(contents)
}

let expected = "e3b0c44298fc1c149afbf4c8996fb92427ae41e4649b934ca495991b7852b855"
let actual = file_checksum("important_data.txt")

if actual == expected {
    print("File integrity verified")
} else {
    print("WARNING: File has been modified!")
}
\end{lstlisting}

\subsection{HMAC: hmac\_sha256}\label{sec:hmac}

\index{hmac\_sha256}

HMAC (Hash-based Message Authentication Code) combines a secret key with a
hash function to verify both the integrity and authenticity of a message:

\begin{lstlisting}[language=Lattice, caption={Computing an HMAC}]
let secret_key = "my-secret-key"
let message = "user=alice&amount=100"

let mac = hmac_sha256(secret_key, message)
print(mac)  // 64-character hex string

// Verify by recomputing
let expected_mac = hmac_sha256(secret_key, message)
if mac == expected_mac {
    print("Message is authentic")
}
\end{lstlisting}

\lstinline|hmac_sha256| takes a key and a message, both as strings, and
returns a hex-encoded MAC. This is the standard way to sign API requests and
verify webhook payloads.

\subsection{Base64 Encoding and Decoding}\label{sec:base64}

\index{base64\_encode}
\index{base64\_decode}

Base64 converts arbitrary binary data into a text-safe ASCII representation:

\begin{lstlisting}[language=Lattice, caption={Base64 encoding and decoding}]
let original = "Hello, World!"

let encoded = base64_encode(original)
print(encoded)  // "SGVsbG8sIFdvcmxkIQ=="

let decoded = base64_decode(encoded)
print(decoded)  // "Hello, World!"
\end{lstlisting}

Base64 is commonly used for embedding binary data in JSON, passing data in
URL parameters, and encoding credentials for HTTP Basic authentication:

\begin{lstlisting}[language=Lattice, caption={HTTP Basic authentication with Base64}]
let username = "admin"
let password = "s3cret"
let credentials = base64_encode(username + ":" + password)

let options = Map::new()
options["headers"] = Map::new()
options["headers"]["Authorization"] = "Basic " + credentials

let response = http_get("https://api.example.com/admin")
\end{lstlisting}

\subsection{Random Bytes}\label{sec:random-bytes}

\index{random\_bytes}

For cryptographically secure random data, use \lstinline|random_bytes|:

\begin{lstlisting}[language=Lattice, caption={Generating random bytes}]
let token_bytes = random_bytes(32)
print(token_bytes.to_hex())  // 64-character random hex string

// Generate a random hex token
fn generate_token(length: Int) -> String {
    let bytes = random_bytes(length)
    return bytes.to_hex()
}

let session_token = generate_token(16)
print("Session: " + session_token)
\end{lstlisting}

\lstinline|random_bytes| returns a \lstinline|Buffer| of the specified
length, filled with cryptographically secure random bytes. On macOS and BSD,
it uses \texttt{arc4random\_buf}. On Linux, it reads from
\texttt{/dev/urandom}. On Windows, it uses \texttt{RtlGenRandom}. With
OpenSSL available, it uses \texttt{RAND\_bytes}. All of these are suitable
for generating session tokens, encryption keys, and nonces.

\begin{defnbox}{Cryptographic vs.\ Pseudorandom}
\lstinline|random_bytes| produces cryptographically secure random data---it
is unpredictable and suitable for security-sensitive applications. The
\lstinline|random()| function from the math module (see \cref{sec:math-random})
uses \texttt{rand()}, which is \emph{not} cryptographically secure and
should only be used for non-security purposes like games and simulations.
\end{defnbox}

% ---------------------------------------------------------------------------
\section{uuid()}\label{sec:uuid}
% ---------------------------------------------------------------------------

\index{uuid}

Universally Unique Identifiers (UUIDs) are 128-bit identifiers typically
displayed as 32 hexadecimal digits in a 8-4-4-4-12 format. Lattice provides
a built-in function for generating Version~4 (random) UUIDs:

\begin{lstlisting}[language=Lattice, caption={Generating UUIDs}]
let id = uuid()
print(id)  // "550e8400-e29b-41d4-a716-446655440000" (example)

// Generate multiple unique identifiers
flux ids = []
flux i = 0
while i < 5 {
    ids.push(uuid())
    i += 1
}
for id in ids {
    print(id)
}
\end{lstlisting}

Each call to \lstinline|uuid()| returns a fresh, unique string. The
underlying implementation uses \lstinline|random_bytes| to generate 16
random bytes, then sets the version bits (0100 for V4) and variant bits
(10xx for RFC~4122) before formatting.

\begin{lstlisting}[language=Lattice, caption={UUIDs for record identifiers}]
fn create_user(name: String, email: String) -> Map {
    let user = Map::new()
    user["id"] = uuid()
    user["name"] = name
    user["email"] = email
    user["created_at"] = time_format(time(), "%Y-%m-%dT%H:%M:%SZ")
    return user
}

let alice = create_user("Alice", "alice@example.com")
print(json_stringify(alice))
\end{lstlisting}

% ---------------------------------------------------------------------------
\section{The os Module}\label{sec:os-module}
% ---------------------------------------------------------------------------

\index{env}
\index{env\_set}
\index{env\_keys}
\index{args}
\index{platform}
\index{cwd}
\index{hostname}
\index{pid}
\index{exec}
\index{shell}

The operating system functions let your program interact with its
environment: reading and setting environment variables, inspecting
command-line arguments, executing external commands, and querying system
information.

\subsection{Environment Variables: env, env\_set, env\_keys}
\label{sec:env-vars}

\begin{lstlisting}[language=Lattice, caption={Working with environment variables}]
// Read an environment variable
let home = env("HOME")
print("Home directory: " + home)

let path = env("PATH")
print("PATH: " + path)

// Check for a variable (returns nil if not set)
let debug = env("DEBUG")
if debug != nil {
    print("Debug mode is on: " + debug)
} else {
    print("Debug mode is off")
}

// Set an environment variable
env_set("APP_MODE", "production")
print(env("APP_MODE"))  // "production"

// List all environment variable names
let keys = env_keys()
print("Environment has " + to_string(keys.len()) + " variables")
\end{lstlisting}

\lstinline|env(name)| returns the value of the named environment variable as
a \lstinline|String|, or \lstinline|nil| if it is not set. \lstinline|env_set|
sets a variable for the current process (it does not modify the parent
shell's environment). \lstinline|env_keys| returns an array of all
environment variable names.

\subsection{Command-Line Arguments: args()}\label{sec:args}

\begin{lstlisting}[language=Lattice, caption={Accessing command-line arguments}]
// Run: lattice script.lat --verbose output.txt

let arguments = args()
print(arguments)  // ["script.lat", "--verbose", "output.txt"]

// Simple argument parsing
flux verbose = false
flux output_path = "default.txt"

flux i = 0
while i < arguments.len() {
    if arguments[i] == "--verbose" {
        verbose = true
    } else if !arguments[i].starts_with("--") {
        output_path = arguments[i]
    }
    i += 1
}

if verbose {
    print("Writing to: " + output_path)
}
\end{lstlisting}

\lstinline|args()| returns the command-line arguments as an array of strings.
The first element is typically the script name.

\subsection{System Information: platform, cwd, hostname, pid}
\label{sec:system-info}

\begin{lstlisting}[language=Lattice, caption={Querying system information}]
print("Platform: " + platform())     // "macos", "linux", "windows", or "wasm"
print("Working dir: " + cwd())       // "/home/user/project"
print("Hostname: " + hostname())     // "workstation"
print("Process ID: " + to_string(pid()))  // 12345
\end{lstlisting}

\lstinline|platform()| returns a short string identifying the operating
system. This is useful for writing cross-platform scripts that need to adjust
their behavior:

\begin{lstlisting}[language=Lattice, caption={Platform-specific behavior}]
let separator = if platform() == "windows" { "\\" } else { "/" }
let config_dir = if platform() == "macos" {
    env("HOME") + "/Library/Application Support/myapp"
} else if platform() == "linux" {
    env("HOME") + "/.config/myapp"
} else {
    env("APPDATA") + "\\myapp"
}

if !file_exists(config_dir) {
    mkdir(config_dir)
}
\end{lstlisting}

\subsection{Executing External Commands: exec and shell}
\label{sec:exec-shell}

\begin{lstlisting}[language=Lattice, caption={Running external commands with exec}]
// exec runs a command and returns its stdout as a string
let output = exec("echo 'Hello from the shell'")
print(output)  // "Hello from the shell\n"

// Capture the output of a system command
let git_status = exec("git status --short")
print(git_status)
\end{lstlisting}

\lstinline|exec| runs a command via \texttt{popen}, captures its standard
output, and returns it as a string. If the command exits with a non-zero
status, \lstinline|exec| raises a runtime error.

For more control---including access to both stdout, stderr, and the exit
code---use \lstinline|shell|:

\begin{lstlisting}[language=Lattice, caption={Running commands with shell for full control}]
let result = shell("ls -la /nonexistent")

print("Exit code: " + to_string(result["exit_code"]))
print("Stdout: " + result["stdout"])
print("Stderr: " + result["stderr"])

// Check if a command succeeded
if result["exit_code"] == 0 {
    print("Success!")
} else {
    print("Failed: " + result["stderr"])
}
\end{lstlisting}

\lstinline|shell| returns a Map with three keys: \lstinline|"exit_code"|
(an \lstinline|Int|), \lstinline|"stdout"| (a \lstinline|String|), and
\lstinline|"stderr"| (a \lstinline|String|). Unlike \lstinline|exec|, it
does \emph{not} raise an error on non-zero exit codes, giving you full
control over error handling.

\begin{warnbox}{Shell Injection}
Both \lstinline|exec| and \lstinline|shell| pass the command string directly
to the system shell (\texttt{/bin/sh} on Unix, \texttt{cmd.exe} on Windows).
\textbf{Never} construct command strings from untrusted user input without
sanitization---this is a classic shell injection vulnerability:
\begin{lstlisting}[language=Lattice]
// DANGEROUS: user_input could contain "; rm -rf /"
let result = exec("echo " + user_input)

// SAFER: validate or escape the input first
\end{lstlisting}
\end{warnbox}

% ---------------------------------------------------------------------------
\section{The math Module}\label{sec:math-module}
% ---------------------------------------------------------------------------

\index{math module}
\index{abs}
\index{floor}
\index{ceil}
\index{round}
\index{sqrt}
\index{pow}
\index{min}
\index{max}
\index{clamp}
\index{lerp}
\index{sin}
\index{cos}
\index{tan}
\index{log}
\index{random}
\index{random\_int}

Lattice's math functions are registered as top-level built-ins. They accept
both \lstinline|Int| and \lstinline|Float| arguments (promoting to
\lstinline|Float| when needed) and cover the standard mathematical operations.

\subsection{Rounding and Absolute Value}\label{sec:math-rounding}

\begin{lstlisting}[language=Lattice, caption={Rounding functions}]
print(abs(-42))       // 42
print(abs(-3.14))     // 3.14

print(floor(3.7))     // 3
print(floor(-2.3))    // -3

print(ceil(3.2))      // 4
print(ceil(-2.7))     // -2

print(round(3.5))     // 4
print(round(2.4))     // 2
\end{lstlisting}

\lstinline|floor|, \lstinline|ceil|, and \lstinline|round| all return
\lstinline|Int| values (the C \texttt{floor}/\texttt{ceil}/\texttt{round}
cast to \texttt{int64\_t}). When given an \lstinline|Int|, they return it
unchanged.

\subsection{Powers, Roots, and Logarithms}\label{sec:math-powers}

\begin{lstlisting}[language=Lattice, caption={Powers and roots}]
print(pow(2, 10))     // 1024 (Int)
print(pow(2, -1))     // 0.5 (Float, because negative exponent)
print(sqrt(144))      // 12.0

print(log(math_e()))   // 1.0 (natural log)
print(log2(1024))      // 10.0
print(log10(1000))     // 3.0

print(exp(1))          // 2.718281828... (e^1)
\end{lstlisting}

\lstinline|pow| is smart about types: when both arguments are non-negative
integers, it returns an \lstinline|Int| (with overflow detection that falls
back to \lstinline|Float|). \lstinline|sqrt| always returns a
\lstinline|Float| and raises a domain error for negative inputs.

\subsection{Min, Max, Clamp, and Lerp}\label{sec:math-minmax}

\begin{lstlisting}[language=Lattice, caption={Clamping and interpolation}]
print(min(3, 7))     // 3
print(max(3, 7))     // 7

// Clamp a value to a range
let health = clamp(120, 0, 100)
print(health)  // 100

let brightness = clamp(-10, 0, 255)
print(brightness)  // 0

// Linear interpolation
let midpoint = lerp(0.0, 100.0, 0.5)
print(midpoint)  // 50.0

let quarter = lerp(10.0, 20.0, 0.25)
print(quarter)  // 12.5
\end{lstlisting}

\lstinline|clamp(value, low, high)| constrains a value to the range
[\textit{low}, \textit{high}]. \lstinline|lerp(a, b, t)| computes the
linear interpolation $a + (b - a) \times t$. Both are essential for game
development, animation, and signal processing.

\subsection{Trigonometry}\label{sec:math-trig}

\begin{lstlisting}[language=Lattice, caption={Trigonometric functions}]
let pi = math_pi()

print(sin(0))           // 0.0
print(sin(pi / 2))      // 1.0
print(cos(0))           // 1.0
print(cos(pi))          // -1.0
print(tan(pi / 4))      // ~1.0

// Inverse trig
print(asin(1.0))        // ~1.5707... (pi/2)
print(acos(0.0))        // ~1.5707... (pi/2)
print(atan(1.0))        // ~0.7853... (pi/4)

// atan2 for angle from coordinates
let angle = atan2(1.0, 1.0)
print(angle)  // ~0.7853... (pi/4, i.e., 45 degrees)
\end{lstlisting}

All trigonometric functions work in radians. \lstinline|math_pi()| returns
$\pi$ and \lstinline|math_e()| returns Euler's number $e$, both as
\lstinline|Float| constants.

Hyperbolic functions are also available: \lstinline|sinh|, \lstinline|cosh|,
and \lstinline|tanh|.

\subsection{Random Numbers}\label{sec:math-random}

\begin{lstlisting}[language=Lattice, caption={Generating random numbers}]
// Random float in [0, 1)
let r = random()
print(r)  // e.g., 0.7234...

// Random integer in a range [low, high] (inclusive)
let dice = random_int(1, 6)
print("You rolled: " + to_string(dice))

// Shuffle an array using random
fn shuffle(arr: Array) -> Array {
    flux result = arr
    flux i = result.len() - 1
    while i > 0 {
        let j = random_int(0, i)
        let temp = result[i]
        result[i] = result[j]
        result[j] = temp
        i -= 1
    }
    return result
}

let deck = [1, 2, 3, 4, 5, 6, 7, 8, 9, 10]
let shuffled = shuffle(deck)
print(shuffled)
\end{lstlisting}

\lstinline|random()| returns a \lstinline|Float| in the range $[0, 1)$.
\lstinline|random_int(low, high)| returns an \lstinline|Int| in the
inclusive range $[\text{low}, \text{high}]$.

\subsection{Other Math Functions}\label{sec:math-other}

\index{sign}
\index{gcd}
\index{lcm}
\index{is\_nan}
\index{is\_inf}

\begin{lstlisting}[language=Lattice, caption={Additional math functions}]
// Sign function
print(sign(-42))   // -1
print(sign(0))     // 0
print(sign(7.5))   // 1.0

// Greatest common divisor and least common multiple
print(gcd(12, 8))   // 4
print(lcm(4, 6))    // 12

// Special value detection
print(is_nan(0.0 / 0.0))   // true
print(is_inf(1.0 / 0.0))   // true
print(is_nan(42.0))         // false
\end{lstlisting}

% ---------------------------------------------------------------------------
\section{The regex Module}\label{sec:regex-module}
% ---------------------------------------------------------------------------

\index{regex}
\index{regex\_match}
\index{regex\_find\_all}
\index{regex\_replace}

Regular expressions let you search, match, and transform text using
patterns. Lattice uses POSIX Extended Regular Expressions (ERE) on Unix
systems.

\subsection{Matching: regex\_match}\label{sec:regex-match}

\begin{lstlisting}[language=Lattice, caption={Testing if a string matches a pattern}]
let email = "alice@example.com"
let pattern = "^[a-zA-Z0-9._%+-]+@[a-zA-Z0-9.-]+\\.[a-zA-Z]{2,}$"

if regex_match(pattern, email) {
    print("Valid email format")
} else {
    print("Invalid email format")
}
\end{lstlisting}

\lstinline|regex_match| returns \lstinline|true| if the pattern matches
anywhere in the string. Use \texttt{\^{}} and \texttt{\$} anchors to require
a full-string match.

\begin{lstlisting}[language=Lattice, caption={Case-insensitive matching with flags}]
// The third argument is a flags string
let is_greeting = regex_match("^hello", "Hello World", "i")
print(is_greeting)  // true
\end{lstlisting}

The optional flags string supports:

\begin{center}
\begin{tabular}{ll}
\toprule
\textbf{Flag} & \textbf{Meaning} \\
\midrule
\texttt{i} & Case-insensitive matching (\texttt{REG\_ICASE}) \\
\texttt{m} & Newline-sensitive matching (\texttt{REG\_NEWLINE}) \\
\bottomrule
\end{tabular}
\end{center}

\subsection{Finding All Matches: regex\_find\_all}\label{sec:regex-find}

\begin{lstlisting}[language=Lattice, caption={Finding all pattern matches in a string}]
let text = "Call 555-1234 or 555-5678 for info"
let phone_pattern = "[0-9]{3}-[0-9]{4}"

let numbers = regex_find_all(phone_pattern, text)
for number in numbers {
    print("Found: " + number)
}
// Found: 555-1234
// Found: 555-5678
\end{lstlisting}

\lstinline|regex_find_all| returns an array of all non-overlapping matches.
If no matches are found, it returns an empty array.

\begin{lstlisting}[language=Lattice, caption={Extracting data from structured text}]
let log = """
[2024-01-15] ERROR: Connection timeout
[2024-01-15] INFO: Server started
[2024-01-16] ERROR: Disk space low
[2024-01-16] WARN: Memory usage high
"""

let errors = regex_find_all("\\[.*\\] ERROR: .*", log)
print("Found " + to_string(errors.len()) + " errors:")
for error in errors {
    print("  " + error)
}
// Found 2 errors:
//   [2024-01-15] ERROR: Connection timeout
//   [2024-01-16] ERROR: Disk space low
\end{lstlisting}

\subsection{Replacing: regex\_replace}\label{sec:regex-replace}

\begin{lstlisting}[language=Lattice, caption={Replacing patterns in text}]
let text = "The quick brown fox jumps over the lazy dog"

// Replace all vowels
let result = regex_replace("[aeiou]", text, "*")
print(result)  // "Th* q**ck br*wn f*x j*mps *v*r th* l*zy d*g"

// Sanitize user input: remove non-alphanumeric characters
let input = "Hello! <script>alert('xss')</script>"
let clean = regex_replace("[^a-zA-Z0-9 ]", input, "")
print(clean)  // "Hello scriptalertxssscript"
\end{lstlisting}

\lstinline|regex_replace| replaces \emph{all} non-overlapping matches of the
pattern with the replacement string. It returns a new string; the original
is unmodified.

\begin{lstlisting}[language=Lattice, caption={Normalizing whitespace}]
let messy = "Too   many     spaces    here"
let clean = regex_replace("  +", messy, " ")
print(clean)  // "Too many spaces here"
\end{lstlisting}

\begin{notebox}{Regex on Windows}
The regex functions use POSIX \texttt{regex.h}, which is not available in
MinGW on Windows. On Windows builds, the regex functions return stub errors.
A future version may bundle PCRE2 for cross-platform support. See
\filename{src/regex\_ops.c} for details.
\end{notebox}

\subsection{Practical Example: Log Parser}\label{sec:regex-example}

Let's combine regex with file I/O and data formats to build a log file
parser:

\begin{lstlisting}[language=Lattice, caption={A log parser using regex}]
fn parse_log(path: String) -> Array {
    let content = read_file(path)
    let lines = content.split("\n")

    flux entries = []
    let pattern = "^\\[([0-9-]+)\\] ([A-Z]+): (.*)$"

    for line in lines {
        if line.len() == 0 { continue }

        let matches = regex_find_all(
            "[0-9]{4}-[0-9]{2}-[0-9]{2}", line)
        let levels = regex_find_all("[A-Z]{3,}", line)

        if matches.len() > 0 && levels.len() > 0 {
            let entry = Map::new()
            entry["date"] = matches[0]
            entry["level"] = levels[0]
            entry["message"] = line
            entries.push(entry)
        }
    }

    return entries
}

let entries = parse_log("application.log")
let json_output = json_stringify(entries)
write_file("parsed_log.json", json_output)
print("Parsed " + to_string(entries.len()) + " log entries")
\end{lstlisting}

% ---------------------------------------------------------------------------
\section*{Exercises}\label{sec:time-crypto-os-exercises}
% ---------------------------------------------------------------------------

\begin{enumerate}
    \item \textbf{Stopwatch.} Write a stopwatch utility with functions
    \lstinline|start()|, \lstinline|lap()|, and \lstinline|stop()| that
    records split times. Print each lap time and the total elapsed time,
    formatted as \texttt{MM:SS.mmm}.

    \item \textbf{Password Hasher.} Write a function that takes a password
    and a salt (generated with \lstinline|random_bytes|), computes
    \lstinline|hmac_sha256(salt, password)|, and returns a string in the
    format \texttt{salt\$hash} (where the salt is hex-encoded). Write a
    companion \lstinline|verify| function that checks a password against a
    stored hash.

    \item \textbf{Cron Expression Evaluator.} Write a function that takes a
    simple cron-like expression (e.g., \texttt{"*/5 * * * *"} for ``every 5
    minutes'') and the current time, and returns the next time the
    expression will fire. Use the \lstinline|time_*| extraction functions
    and date arithmetic.

    \item \textbf{System Information Report.} Write a program that collects
    system information---platform, hostname, PID, current directory, all
    environment variables, and the current time---and outputs it as a
    formatted YAML document.

    \item \textbf{Text Statistics.} Write a program that reads a text file
    and uses \lstinline|regex_find_all| to count: (a)~the number of words,
    (b)~the number of sentences (ending with . ! or ?), (c)~the number of
    email addresses, and (d)~the number of URLs. Output the results as a
    JSON object.

    \item \textbf{Monte Carlo Pi Estimator.} Use \lstinline|random()| to
    estimate $\pi$ by the Monte Carlo method: generate random $(x, y)$
    points in the unit square, count how many fall inside the unit circle
    ($x^2 + y^2 \leq 1$), and compute $\pi \approx 4 \times
    \text{inside} / \text{total}$. How many points do you need for 4
    decimal places of accuracy?
\end{enumerate}

% ---------------------------------------------------------------------------
\section*{What's Next}\label{sec:time-crypto-os-next}
% ---------------------------------------------------------------------------

With this chapter, we have completed our tour of the standard library.
You now have the tools to work with files, parse and generate data in
multiple formats, communicate over the network, manipulate dates, hash and
encode data, generate unique identifiers, interact with the operating system,
do mathematics, and match text with regular expressions. In
\cref{ch:lazy-iterators}, we will shift gears and explore Lattice's iterator
system---a powerful abstraction for processing sequences lazily and
composing transformations with a functional style.
